\chapter{Command-Line Interface}
\label{ch:cli}

\newthought{Everything you do with \Jarvis} goes through one command, \cmd{Jarvis}. This chapter is
the reference for its options---running a scan, converting and plotting results, monitoring a live
run, and checking a card before you commit to it.

\section{Two entry points}
\label{sec:cli-entry}

\cmd{Jarvis} has exactly two forms:

\begin{shellbox}
Jarvis <file.yaml> [options]      # run (or act on) a scan
Jarvis project <command> [args]   # manage standalone projects
\end{shellbox}

\noindent The first runs a task card; the second manages projects and is covered in
Chapter~\ref{ch:project-cli}. This chapter is about the first.

\section{Running a scan}
\label{sec:cli-run}

With no options, \cmd{Jarvis} runs the scan described by the card:

\begin{shellbox}
Jarvis bin/my_scan.yaml
\end{shellbox}

\section{The options}
\label{sec:cli-options}

Table~\ref{tab:cli} lists every option. They fall into general flags, workflow utilities that act
on a run, and run-mode switches.

\begin{table}[ht]
  \footnotesize
  \begin{tabular}{@{}lp{0.60\linewidth}@{}}
    \toprule
    \textbf{Option} & \textbf{What it does} \\
    \midrule
    \cli{-h}, \cli{-{}-help}    & show usage and exit \\
    \cli{-v}, \cli{-{}-version} & print the banner and installed version \\
    \cli{-{}-refs}             & print full reference information \\
    \cli{-d}, \cli{-{}-debug}   & lower the log level to \code{debug} (verbose) \\
    \cli{-{}-resume}           & resume from the latest checkpoint without prompting (Ch.~\ref{ch:checkpoint}) \\
    \cli{-{}-convert}          & convert the live \textsc{hdf5} database to \textsc{csv} \\
    \cli{-{}-plot}             & generate a \JPlot\ configuration from the card \\
    \cli{-{}-monitor}          & open a real-time resource monitor for a running scan \\
    \cli{-{}-check-modules}    & test-run a few samples to validate calculators/modules \\
    \cli{-{}-benchmark [s]}    & run a throughput benchmark (default 30\,s) \\
    \cli{-{}-skip-library-installation} & skip the \yamlkey{LibDeps} install step \\
    \cli{-{}-skip-draw-flowchart} & skip the \textsc{png} flowchart (still writes \file{flowchart.json}) \\
    \bottomrule
  \end{tabular}
  \caption{The \cmd{Jarvis} command-line options.}
  \label{tab:cli}
\end{table}

\section{Utilities that act on a run}
\label{sec:cli-utilities}

Three options operate on an \emph{existing} run rather than starting a fresh scan---typically in a
second terminal while a long scan is going:

\begin{description}[style=nextline, leftmargin=1.2em]
  \item[\cli{-{}-convert}] Snapshot the current \textsc{hdf5} database into \textsc{csv} so you can
        inspect partial results (or feed \JPlot) without stopping the scan.
  \item[\cli{-{}-plot}] Write a \JPlot\ \textsc{yaml} for the run; render it with \cmd{jplot}
        (Chapter~\ref{ch:jarvisplot}).
  \item[\cli{-{}-monitor}] Open a live resource monitor that refreshes every few seconds; press
        \code{q} to quit.
\end{description}

\begin{jwarn}
\cli{-{}-convert} and \cli{-{}-monitor} work on a \emph{live or finished} run's data. In
particular, \cli{-{}-monitor} attaches to the running process, so start it while the scan is
running---the process is gone once the scan ends.
\end{jwarn}

\section{Checking a card before a full run}
\label{sec:cli-check}

\cli{-{}-check-modules} runs a small test scan (a handful of samples) so you can confirm your
calculators install, run, and produce the expected observables---before launching millions of
points:

\begin{shellbox}
Jarvis bin/my_scan.yaml --check-modules
\end{shellbox}

\noindent It writes its per-sample artifacts under \file{outputs/<scan>/SAMPLE/tests/} (the
\file{DATABASE/} path is shared with a normal run), so a check never pollutes real results.

\section{Measuring throughput}
\label{sec:cli-benchmark}

\cli{-{}-benchmark} runs the machinery in a throughput mode for a fixed number of seconds (default
30) and reports points-per-second---useful for tuning workers and comparing machines
(Chapter~\ref{ch:performance}):

\begin{shellbox}
Jarvis bin/my_scan.yaml --benchmark 60
\end{shellbox}

\noindent As well as printing the rate, it writes a \file{benchmark.json} with the measured
\code{samples\_per\_sec}, the wall time, and the submitted/completed/failed counts---a reproducible
record you can compare across machines or settings.

\section{Where to go next}
\label{sec:cli-next}

\begin{itemize}
  \item Project scaffolding and packaging: Chapter~\ref{ch:project-cli}.
  \item Resuming an interrupted scan: Chapter~\ref{ch:checkpoint}.
  \item What the run produced: Chapter~\ref{ch:outputs}.
\end{itemize}
